Every FORTRAN file that references {\CDI} functions or constants must contain
an appropriate {\tt INCLUDE} statement before the first such reference:

\begin{verbatim}
   INCLUDE "cdi.inc"
\end{verbatim}

Unless the {\tt cdi.inc} file is installed in a standard directory where
FORTRAN compiler always looks, you must use the {\tt -I} option when invoking
the compiler, to specify a directory where {\tt cdi.inc} is installed, for example:

\begin{verbatim}
   f77 -c -I/usr/local/cdi/include myprogram.f
\end{verbatim}

Alternatively, you could specify an absolute path name in the {\tt INCLUDE}
statement, but then your program would not compile on another platform
where {\CDI} is installed in a different location.

Unless the {\CDI} library is installed in a standard directory where the linker
always looks, you must use the {\tt -L} and {\tt -l} options to links an object file that
uses the {\CDI} library. For example:

\begin{verbatim}
   f77 -o myprogram myprogram.o -L/usr/local/cdi/lib -lcdi
\end{verbatim}

Alternatively, you could specify an absolute path name for the library:

\begin{verbatim}
   f77 -o myprogram myprogram.o -L/usr/local/cdi/lib/libcdi
\end{verbatim}

If the {\CDI} library is using other external libraries, you must add this
libraries in the same way.
For example with the netCDF library:

\begin{verbatim}
   f77 -o myprogram myprogram.o -L/usr/local/cdi/lib -lcdi \
                                -L/usr/local/netcdf/lib -lnetcdf
\end{verbatim}
